% Generated draft: requires mathematical and source review.
\documentclass[11pt]{article}
\usepackage[margin=25mm]{geometry}
\usepackage{fontspec}
\setmainfont{Latin Modern Roman}
\usepackage{amsmath,amssymb,mathtools,unicode-math}
\setmathfont{Latin Modern Math}
\usepackage{enumitem,needspace,xurl,hyperref}
\setlength{\parindent}{0pt}
\setlength{\parskip}{5pt}
\setlength{\emergencystretch}{3em}
\newcommand{\N}{\mathbb{N}}
\newcommand{\Z}{\mathbb{Z}}
\newcommand{\Q}{\mathbb{Q}}
\newcommand{\R}{\mathbb{R}}
\newcommand{\C}{\mathbb{C}}
\newcommand{\F}{\mathbb{F}}
\newcommand{\A}{\mathbb{A}}
\newcommand{\PP}{\mathbb{P}}
\newcommand{\E}{\mathbb{E}}
\newcommand{\eps}{\varepsilon}
\newcommand{\dd}{\,\mathrm{d}}
\begin{document}
\section*{UnsolvedMath ID: 701301}
\textbf{Source problem code:} catalogue entry
\subsection{Definitions and mathematical statement}
Let \(F_n\) be a free group of finite rank \(n\), equipped with a fixed free basis \(X=\{x_1,\ldots,x_n\}\). Elements are represented by finite words in \(X\cup X^{-1}\). For \(w\in F_n\), let \(|w|_X\) denote its freely reduced word length, and define its cyclic length by
\[
\|w\|_{\mathrm{cyc},X}=\min_{g\in F_n}|g^{-1}wg|_X.
\]
Equivalently, this is the length of a cyclically reduced representative of the conjugacy class of \(w\); the identity has cyclic length zero. Write \(\operatorname{Aut}(F_n)\) for the group of automorphisms of \(F_n\).

The question is whether there is an algorithm which, given arbitrary \(u,v\in F_n\), terminates and decides whether
\[
\forall\phi\in\operatorname{Aut}(F_n),\qquad
\|\phi(u)\|_{\mathrm{cyc},X}=\|\phi(v)\|_{\mathrm{cyc},X}.
\]
The source does not explicitly specify the range of ranks \(n\), or whether it asks for an algorithm for each fixed rank or one algorithm uniform in \(n\). These algorithmic scope qualifications remain unresolved.

The supplied Comment of 2009 reports an affirmative answer for \(n=2\), citing D. Lee, \emph{J. Group Theory} \textbf{10} (2007), 561--569. This is a supplied historical qualification. The registry labels the overall problem open; neither that label nor the reported rank-two result has been independently verified here.
\subsection{Short English statement}
Given two elements of a free group, decide whether every automorphism sends them to elements with equal cyclically reduced lengths. The supplied source reports that an algorithm exists in rank two. The intended rank range and whether the algorithm must be uniform in the rank are unspecified.
\subsection{Sources}
\begin{itemize}
\item Local UnsolvedMath registry, permanent ID \(701301\), exact source code \texttt{catalogue entry}, title \emph{Kourovka Notebook Problem 16.110}; supplied upstream revision \texttt{d371ee917ac9647fbb06cd6fc529deec6f3dda82}. The supplied registry status is \texttt{open}, not an independently verified literature assessment.
\item \emph{The Kourovka Notebook (Unsolved Problems in Group Theory)}, Problem 16.110, 16th issue (2006); the background also identifies the 21st issue edition (2026), \texttt{arXiv:1401.0300v48}. The supplied collection description names E. I. Khukhro and V. D. Mazurov as editors.
\item Attribution as supplied: the statement begins with I. Kapovich and P. Schupp in parentheses, while the proposer field and background name V. E. Shpilrain, with proposal year 2006. The relationship between these attributions is not explained.
\item Supplied Comment of 2009: affirmative answer for \(n=2\), attributed to D. Lee, \emph{J. Group Theory} \textbf{10} (2007), 561--569. No article title or direct link is supplied.
\end{itemize}
\end{document}
